\chapter{Estabilidad Espectral Asintótica, Factorización $LDL^T$ de Rook y QSBR (V815--V816)}
\label{cap:v815_estabilidad_espectral}

\section{Introducción y Diagnóstico del Tribunal Multi-IA}
En la iteración V815/V816 del Kernel POLYDIM, el Tribunal Adversarial de IAs (Cerebras CS-3 120B, Kimi Moonshot, DeepSeek Coder y Claude Sonnet) expuso límites asintóticos fundamentales en la reducción de Schur para el operador de retracción de Cayley-SMW bajo regímenes de alto número de condición, así como la necesidad de blindaje en la gestión de memoria lock-free para topologías multi-agente en silicio heterogéneo.

\section{Condicionamiento de la Reducción de Schur en Cayley-SMW}
Consideremos la retracción ortogonal sobre la variedad de Stiefel $St(D, K)$ calculada mediante la reducción de Schur $2K \times 2K \to K \times K$:
\begin{equation}
M = I_K + \alpha (S - S^T) + \alpha^2 Q, \quad S \in \mathbb{R}^{K \times K}, \; Q = S S^T
\end{equation}
Sean $\sigma_i = \sigma_i(S - S^T)$ los valores singulares de la componente antisimétrica. Los autovalores de $M$ satisfacen:
\begin{equation}
\lambda_i(M) = 1 + \alpha \sigma_i + \alpha^2 \lambda_i(Q)
\end{equation}

\begin{theorem}[Cota de Exponentiación y Divergencia de Condición]
Si $|\alpha| \sigma_{\max}(S - S^T) \gg 1$, el autovalor mínimo de $M$ decae asintóticamente como:
\begin{equation}
\min_i |\lambda_i(M)| = O(|\alpha|^{-1})
\end{equation}
lo que induce un número de condición $\kappa(M) = \frac{\max_i |\lambda_i|}{\min_i |\lambda_i|} = O(|\alpha| \sigma_{\max})$.
\end{theorem}

\begin{proof}
Para matrices antisimétricas genéricas, la componente imaginaria de los autovalores desvía la trayectoria del espectro en el plano complejo. Cuando $\alpha$ no está acotado respecto al radio espectral de $S - S^T$, la aproximación de primer orden se cancela con la identidad en direcciones degeneradas, colapsando $\min |\lambda_i| \to 0$ y provocando inestabilidad numérica en precisión IEEE-754 doble ($\kappa(M) \ge 10^{11}$).
\end{proof}

\section{Teorema de Normalización y Scaling Espectral}
Para garantizar la exactitud de la retracción sin pérdida de ortogonalidad, se establece la ley de scaling espectral dinámico:

\begin{proposition}[Spectral Scaling Invariante]
Definiendo el factor de paso normalizado:
\begin{equation}
\alpha^* = \frac{\alpha}{\max\left(1, \; |\alpha| \cdot \sigma_{\max}(S - S^T)\right)}
\end{equation}
se garantiza que $\kappa(M) \le 1 + 2|\alpha^*| \sigma_{\max} + (\alpha^*)^2 \|Q\|_2 \le O(1)$, asegurando la estabilidad numérica incondicional del resolvedor lineal.
\end{proposition}

\section{Factorización $LDL^T$ por Bloques con Pivoteo de Rook}
Para mitigar la degeneración inducida por pivotes nulos en matrices simétricas indefinidas derivadas de proyecciones tangenciales, el Kernel V816 implementa la descomposición $P M P^T = L D L^T$ mediante pivoteo de Rook acotado:
\begin{equation}
D = \operatorname{diag}(D_1, D_2, \dots, D_m), \quad D_j \in \{\mathbb{R}^{1 \times 1}, \mathbb{R}^{2 \times 2}\}
\end{equation}
El pivoteo de Rook garantiza que los elementos de $L$ satisfagan $|L_{ij}| \le \frac{1}{1 - \mu}$ con costo $O(K^2)$ por búsqueda, evitando el costo $O(K^3)$ del pivoteo completo de Bunch-Kaufman sin comprometer la estabilidad hacia atrás.

\section{Reclamación de Memoria Basada en Estados Quiescentes (QSBR Arena)}
En el protocolo PMTP de cero copias, para evitar el problema ABA y la degradación por contención de bloqueos atómicos en arquitecturas NUMA multi-socket, se introduce el gestor de memoria QSBR:
\begin{enumerate}
    \item Cada hilo lector publica periódicamente un epoch monotónico en una línea de caché privada aislada (128 bytes padding anti false-sharing).
    \item La transición de slabs liberados sigue el ciclo estricto:
    \begin{equation}
    \text{ACTIVE} \xrightarrow{\text{CAS}} \text{SUSPECT} \xrightarrow{\text{Acquire Fence}} \text{QUIESCENT} \xrightarrow{\text{Reclaim}} \text{FREE}
    \end{equation}
    \item Ningún buffer es reasignado hasta que $\min_{t} (\text{Epoch}_t) > \text{Epoch}_{\text{retire}}$, garantizando ausencia total de Use-After-Free (UAF) sin penalización de locks en el camino crítico.
\end{enumerate}

\section{Preservación de Fase Cuántica en el Isomorfismo $Sp(2n, \mathbb{F}_2)$}
El compilador del puente de Clifford+T recompone el grupo simpléctico preservando la fase global:
\begin{equation}
U = U_{\text{Clifford}} \cdot \prod_{k=1}^{n} T_k \cdot D_{\text{phase}}(\theta), \quad \theta \in [0, 2\pi)
\end{equation}
impidiendo la decoherencia de fase en codiagonalizaciones de alta dimensión ($n \ge 4$).

\section{Refutación de Truncamiento en BF16 y Prevención de Cancelación en Origen}
\label{sec:refutacion_bf16_newton_schulz}

Existe una falacia recurrente en la literatura de bajo costo computacional: asumir que es posible calcular gradientes geodésicos en precisión reducida $\text{bfloat16}$ ($\epsilon_{\text{BF16}} \approx 7.81 \times 10^{-3}$) y posteriormente ``recuperar'' la geometría ortogonal mediante iteraciones de Newton-Schulz en doble precisión ($\text{FP64}$).

\begin{theorem}[Irrecuperabilidad Informacional por Cancelación Catastrófica]
Sea $\mathbf{g} \in T_{\mathbf{x}}\mathcal{M}$ en dimensión $D \ge 10^6$. Si la proyección al formato $\text{fl}_{\text{BF16}}(\mathbf{g})$ induce cancelación catastrófica de componentes ortogonales subordinadas ($|g_i - g_j| < 2^{-8} \max(|g_i|, |g_j|)$), la pérdida de información es termodinámicamente irreversible por el Teorema de Procesamiento de Datos (DPI). La subsecuente aplicación de Newton-Schulz en $\text{FP64}$:
\begin{equation}
X_{k+1} = X_k \left( \frac{3}{2} I - \frac{1}{2} X_k^T X_k \right)
\end{equation}
únicamente ortogonaliza una matriz degenerada $\tilde{X} = X + E_{\text{BF16}}$, preservando el ruido numérico y acelerando el drift tangencial fuera de la geodésica.
\end{theorem}

\subsection{Soluciones SOTA 2025--2026: Prevención en Origen}
Para neutralizar la cancelación catastrófica en $S^{D-1}$, el Kernel POLYDIM adopta los siguientes mecanismos certificados:
\begin{enumerate}
    \item \textbf{Precondicionamiento de Segundo Momento en FP32 (Paradigma Muon$^2$):} Se acumulan los momentos adaptativos estrictamente en $\text{FP32}$ ($\tilde{B}_t = B_t / (\sqrt{V_t} + \epsilon)$) antes de cualquier paso de retracción, comprimiendo el espectro de valores singulares fuera de la zona de singularidad.
    \item \textbf{Adaptación por Dimensión (AdaNewton / ROOT):} Se eliminan los coeficientes fijos de Newton-Schulz, sustituyéndolos por tablas pre-optimizadas $\{a^{(m,n)}, b^{(m,n)}, c^{(m,n)}\}$ por dimensión $(m,n)$ y filtrado \textit{soft-threshold} contra colas pesadas.
    \item \textbf{Acotamiento de Iteraciones (Dao Lab Bound):} Se limita rigurosamente $q \le 2$ iteraciones en aproximaciones de Newton-Schulz para impedir la generación de autovalores negativos espurios en $Q_t = X_t^T X_t$.
    \item \textbf{Región de Confianza Espectral Lineal (AuON):} Para $D \ge 10^6$, se sustituye la ortogonalización $O(n^2)$ por normalización espectral y escalado no lineal acotado en $O(n)$.
\end{enumerate}

